% ----------------------------------------------------------------------
% Chapter 5: Conclusion
% Standalone compilable artifact (own preamble); when assembling the
% thesis, strip the preamble and merge the body only.
% Short: summary plus a brief outlook, per ch5_structure.md.
% Own results are referred to by name ("the attainment theorem of
% Chapter 3", "the approximation theorem of Chapter 4", ...) rather
% than by draft numbers; FAKE REF comments mark where real \ref's
% should be inserted on final merge.
% ----------------------------------------------------------------------
\documentclass[11pt,a4paper]{report}
\usepackage[T1]{fontenc}
\usepackage[utf8]{inputenc}
\usepackage{amsmath,amssymb,amsthm}
\usepackage{enumitem}

\numberwithin{equation}{chapter}

\newcommand{\E}{\mathbb{E}}
\newcommand{\PP}{\mathbb{P}}
\newcommand{\R}{\mathbb{R}}
\newcommand{\F}{\mathcal{F}}
\newcommand{\HH}{\mathcal{H}}

\begin{document}

\setcounter{chapter}{4}

\chapter{Conclusion}

\section{Summary}

This thesis develops the deep hedging framework of B\"uhler, Gonon, Teichmann and Wood \cite{buehler} in finite discrete time on a general probability space. Chapters~1 and~2 set up a market with trading constraints and transaction costs, with the terminal profit-and-loss evaluated by an optimized certainty equivalent (OCE). Already here, domain and integrability of the risk measure matter, and additional assumptions are needed to make the risk measure sensible.

Chapter~3 defines addiontal assumptions: (A1)-(A6) (bounded liability, convex strategy set, half-bounded Fatou property, $L^0$-boundedness of the gain set, closedness of the dominated gain cone and nonnegative gains). We then proof that they repair the existence of an optimal hedge, which would otherwise fail on a general probability space. Chapter~4 treats neural network approximation. On finite $\Omega$, $\pi_M\to\pi$, neural networks approximate the optimal hedge but fail on a general probability space. Again with (A1) and (A6) we restore $\pi_M(-Z)\to\pi(-Z)$ for OCEs. Thus under (A1)-(A6) network strategies are $\varepsilon$-optimal for an \emph{attained} optimum.

Two natural expansions remain. Extending the approximation theory from OCEs to general convex risk measures via robust representation \cite[Section~4.5]{buehler} is open on general $\Omega$. Furthermore, I claimed (A1)-(A6) to be economically sensible. Whether they hold, at least approximately, in real markets is still an empirical question that needs to be investigated.

\begin{thebibliography}{9}

    \bibitem{buehler}
    H.~B\"uhler, L.~Gonon, J.~Teichmann, and B.~Wood.
    Deep hedging.
    \emph{Quantitative Finance} \textbf{19}(8), 1271--1291, 2019.

    \bibitem{delbaenschachermayer}
    F.~Delbaen and W.~Schachermayer.
    A general version of the fundamental theorem of asset pricing.
    \emph{Mathematische Annalen} \textbf{300}(1), 463--520, 1994.

\end{thebibliography}

\end{document}
